\documentclass[10pt,letterpaper]{article}
\usepackage{amsmath}
\usepackage{amsfonts}
\usepackage{amssymb}
\usepackage[english]{babel}
\usepackage{breakurl}
\usepackage[superscript]{cite}
\usepackage{fancyhdr}
\usepackage{float}
\usepackage[margin=1in]{geometry}
\usepackage{graphicx}
\usepackage{hyperref}
\usepackage[utf8]{inputenc}
\usepackage{makeidx}
\usepackage{multicol}
\usepackage{nth}
\usepackage{longtable}

% Use Helvetica
\usepackage[scaled]{helvet}
\renewcommand\familydefault{\sfdefault} 
\usepackage[T1]{fontenc}

\usepackage{setspace}
\usepackage{siunitx}
\usepackage{svg}
\usepackage{subcaption}
\usepackage{tikz}
\usepackage{titling}
\usepackage{url}
\usepackage{kantlipsum}
\usepackage{enumerate}

\usepackage{parskip} % Paragraph skip - adds extra lineskip spacing
\setlength{\parskip}{0.7\baselineskip plus 2pt}

% Custom definitions
\newcommand{\doctitle}{Requirements Specification}
\newcommand{\docsubtitle}{Computing Platform Multirotor with FPGA Hardware Acceleration Applications}

% Custom commands
\newcommand{\ts}{\textsubscript}	% Subscript command %

% Use hyphans to break up urls
\def\UrlBreaks{\do\/\do-}

% PDF and href setup
% Hyper ref
\hypersetup{
	colorlinks=true,
	citecolor=black,
	linkcolor=black,
	filecolor=black,
	urlcolor=blue,
	pdftitle={\@title},
	bookmarks=true
}
\urlstyle{same}

% Page headings
\pagestyle{fancy}
\fancyhead[L]{\MakeUppercase{CPEN/ELEC 491}}
\fancyhead[R]{\textbf{Team 109}}
\fancyfoot{}
\fancyfoot[C]{\thepage}

% No paragraph indent
\parindent 0ex

% Meta
\author{
	Deutsch, Peter &
	\textit{me@peterdeutsch.ca}
	\\
	He, Muchen &
	\textit{i@muchen.ca}
	\\
	Hsueh, Arthur &
	\textit{ah11962@outlook.com}
	\\
	Wang, Meng &
	\textit{wzfftxwd@gmail.com}
	\\
	Wilson, Ardell &
	\textit{ardellw96@gmail.com}
}
\title{\doctitle}
\date{\today}

\makeatletter
\renewcommand{\maketitle}{
	\bgroup
	\setlength{\parindent}{0pt}
	\begin{flushleft}
		% Top spacing
		\vspace*{0.75in}

		% Team logo
		\includegraphics[scale=0.5]{../assets/capstonelogo1.png}
		\vspace*{0.25in}

		% Title
		\textbf{\Huge{\@title}}\\
		\hrulefill

		% Subtitle
		\textbf{\huge{\docsubtitle}}
		
		\vspace*{0.5in}

		% Course number and team
		\textbf{\Large{CPEN/ELEC 491 Capstone Team 109}}\\
		\hspace*{0.1cm}
		\begin{tabular}[h]{|ll}
			\@author
		\end{tabular}

		\vspace*{0.25in}

		\textbf{Website}: https://capstone-skynet.github.io

		\vfill

		% Logo
		\hspace*{-0.3cm}\includegraphics[scale=0.5]{../assets/ece_logo.pdf}

		% Date
		\large{\@date}
	\end{flushleft}
	\egroup
}
\makeatother

% Begin Document
\begin{document}

% Title Page
\begin{titlepage}
	\maketitle
\end{titlepage}

% Revision history
%\thispagestyle{empty}
\section*{Revision History}
The full revision history and commited changes of the document can be found in the git repository history: \href{https://github.com/Capstone-Skynet/Capstone-Skynet.github.io}{https://github.com/Capstone-Skynet/Capstone-Skynet.github.io/commits/master}.

\begin{table}[H]
\begin{tabular}{*{4}{l}p{0.5\linewidth}}
\hline
Version \# & Initials & Release Date & Changeset & Changes Made \\ \hline

0.0 & PD & 2019-10-11 & \texttt{660e001} & Initial skeleton of the document.\\
0.1 & MH & 2019-10-11 & \texttt{6af9e8a} & Populate initial document with draft content required for Milestone I.\\
0.2 & PD & 2019-11-23 & & Initial framework for test descriptions created.\\
0.3 & MW & 2019-11-24 & & First set of tests added.\\
1.0 & PD & 2019-11-25 & & General clean-up and release for Milestone II.\\
1.1 & AH & 2020-02-09 & & Added Full System Tests \\

 & & & \\ \hline
\end{tabular}
\end{table}

% Table of contents
\setcounter{secnumdepth}{3}
\tableofcontents

% Terms and Abbreviations
\thispagestyle{empty}

\section*{Terms and Abbreviations}

Technical terms and abbreviations dictionary go here.

\begin{tabular}[h]{rp{0.75\linewidth}}
    \hline
    \textbf{Term} & \textbf{Definition}\\
    \hline

    ANN & Artificial Neural Network, or simply ``Neural Network'', is a data processing model modeled after neuron interactions. The process consists of forward propagation using several matrix multiplications.\cite{ann}\\
    ASIC & Application-specific Integrated Circuits.\\
    CNN & Convolutional neural network are neural networks that is especially useful for image classification.\cite{cnn} \\
    ECE & Department of Electrical and Computer Engineering at the University of British Columbia.\\
    FPGA & Field-Programmable-Gate-Arrays, ``programmable'' hardware that allows ASIC-like performance with software-like turn-around time and flexibility.\\
    GPU & Graphics Processing Unit, a discrete piece of hardware designed to accelerate graphic-intensive or other parallel computing tasks.\\
    LOS & Line-of-sight.\\
    ML & Machine learning.\\
    Multirotor & An unmanned vehicle with multiple engines. \\
    OTS & Off-the-shelf, or commercially available/purchasable \\
    PID / PID Controller & Proportion-Integral-Derivative controllers is the most common control algorithm for precise and accurate movement, as well as to compensate external forces.\cite{pid}\\
    RNN & Recurrent neural networks are neural networks where the output depends on previous computations, essentially consists of memory.\cite{rnn}\\
    RX & Receiver.\\
    TC & Transport Canada.\\
    TX & Transmitter.\\
    YOLO & You-Only-Look-Once is a fast ML algorithm that detect objects but is unlike CNN nor RNN.\cite{yolo}\cite{yolo-2}\\
     & \\

    \hline

\end{tabular}


% List of figures and tables
\thispagestyle{empty}
\listoffigures
\listoftables
\newpage

% Set page and section counter
\setcounter{page}{1}

% Begin content
\section{About This Document}\label{section:about}
\subsection{Purpose}\label{section:about:purpose}
This document outlines the project's requirements, laying out a basis of what is required for the project to be deemed \textit{successful}. In particular, this document outlines expected metrics related to performance, quality/longevity, environmental impact, cost, legality, and safety.

\subsection{Intended Audiences}\label{section:about:audience}
This document is intended for \textit{all project stakeholders}, including (but not limited to):
\begin{itemize}
\item Project Designers and Quality Controllers (the Capstone team)
\item Client (and their representatives)
\item External Project Sponsors
\item Legal Personnel (The Department of Electrical and Computer Engineering, Industry Canada, Transport Canada)
\end{itemize}

\subsection{Reading Guide}\label{section:about:readingguide}
This document describes both a high-level description of the project, in addition to detailed functional/non-functional requirements. Each detailed requirement listed in this document is assigned a unique \textit{requirement identifier}, which can be cross-referenced to other project documents (mainly the Design Specification). 

\section{Background}

% 
% Describes the importance of a flying FPGA ML thing to the client (Mieszko).
% 

TODO: A nice high-level pictogram (\textit{not} a block diagram, that's for the design doc) would be great here.

\section{Domain}

The project is classified as consumer electronics integrated with components used in research and industry such as FPGA hardware accelerators. 

\section{Goal}

This project is different than other products on the market as we offer programmable hardware that could be taylored to a wide range of applications that require hardware acceleration, prodominately machine learning models. Popular products on the market are mostly aimed at the consumers and ``prosumers'' or content-producers with a main focus in aerial photography. While our project also features aerial photography, the emphasis is in performing real-time analytics using the captured video as well as custom features and modifications taylored to our client.

\section{Functional Requirements}\label{section:funcreq}

This section outlines the functional requirements of the integreated system.

Functional requirements describes what the project do. It embodies the essential requirements that must be completed in order for the project to be deemed deliverable or successful. The functional requirements include the scope of business rules, legal or regulartory requirements, certification requirements, and external interfacing requirements for the project.

\subsection{Multirotor RPAS}
\begin{enumerate}[{F.DR}.1:]
	\item Total takeoff mass of the drone with payload is less than 25kg
	\item Drone is self-stabilizing using its onboard flight controller.
	\item Drone is capable of being controlled by a pilot within line-of-sight.

\end{enumerate}

\subsection{Non-RPAS Air-Ground I/O}
\begin{enumerate}[{F.CM}.1:]
	\item Computation platform transmits video data and ML model data to ground station.
\end{enumerate}

\subsection{Computing Platform}
\begin{enumerate}[{F.CP}.1:]
	\item Captures video at a minimum resolution of 640x480 at 10fps.
	\item Onboard FPGA is capable of performing parallel matrix computation to speedup machine learning.
\end{enumerate}

% Camera and ML points
\subsection{Camera}
\begin{enumerate}[{F.CAM}.1]
	\item The camera must capture video for the machine learning implementation to process
\end{enumerate}

\subsection{Machine Learning Implementation}
\begin{enumerate}[{F.ML}.1:]
	\item A machine learning model, implemented on the FPGA, is capable of detecting pedestrians from image frames. captured using the camera.
	\item The machine learning model must be supported by FPGA-based hardware acceleration.
	\item Any Intellectual Property (IP) cores used by the hardware for machine learning must be compatible with the MIT license.
\end{enumerate}

\subsubsection{Base Station}
\begin{enumerate}[{F.BS}.1:]
    \item Ground station displays the video data overlaid with the ML model data.
	\item ML model data is presented in the form of bounding boxes.
\end{enumerate}

\section{Non-Functional Requirements}\label{section:nonfuncrec}

TODO: We might have to split non-functional requirements into nf-reqs and constraints (the capstone docs say they're separate, but PC's don't).

% Camera and ML points
\subsection{Camera}
\begin{enumerate}[{NF.CAM}.1]
	\item The cameral must use a standard interface, for future replacability
	\item The camera must be capable of focusing on objects at variable distances, between 1 and 30 meters.
	\item The camera must resolve individual objects from a minimum height of 10 meters to a precision of 10cm across.
	\item The camera must support video capture at a minimum rate of 10 frames per second.
	\item The camera must support a minimum resolution of 640x480.
\end{enumerate}

\subsection{Machine Learning Implementation}
\begin{enumerate}[{NF.ML}.1]
	\item The machine learning model must be implemented for novel video processing applications.
	\item The machine learning model must have a throughput of at least 2 frames per second.
	\item The machine learning model must process in real time, with a latency of one second or less.
	\item The machine learning model must process both live inputs (from the camera) and test inputs (from a pre-recorded source). 
	\item The machine learning model must output metrics based on objects within a given video frame
\end{enumerate}

\subsection{Multirotor RPAS}
\begin{enumerate}[{NF.FP}.1:]
    \item Drone has at least 50\% flight duration of that without the payload --- 10 minutes minimum.
    \item The drone must be able to fly in fair conditions (assuming no rain, heavy winds, etc.).
    \item The drone should be capable of auto-leveling, drift-correction, and perform emergency landing.
    \item The drone should have systems for low-battery alarm and emergency protocol for low-battery and loss of control radio.
\end{enumerate}

\subsection{System Reliability and Durability}
\begin{enumerate}[{NF.SR}.1:]
    \item Brushless DC Motors used on the drone should have a life span of 10,000 flight-hours.
    \item Non-computation-platform electronics are expected to have a life span of 20,000 flight hours if the procedures are followed correctly
    \item The lithium polymer battery for the drone is expected to retain 75\% of its charge after 200 charge/discharge cycles. (TODO: make citation \url{http://learningrc.com/lipo-battery/}) 
\end{enumerate}

\section{Contraints}

Constraints are not part of functional or non-function specifications.

\subsection{Cost}
\begin{enumerate}[{C.CT}.1:]
    \item Total cost of the project should be below C\$1,000, with C\$650 provided by the Capstone course as a backup.
    \item Operating cost of the integrated drone system upon completion should be practically free (non-including electricity cost, cost of time and labour, opportunity cost of the client, and cost for transportation).
\end{enumerate}

\subsection{Platform Extensibility}
\begin{enumerate}[{C.PE}.1:]
    \item Significant on-chip area (>50\%) should be dedicated for future client use.
\end{enumerate}

\subsection{Legal Compliance}
\begin{enumerate}[{C.LG}.1:]
    \item Drone is in compliance with Transport Canada's new laws (TC 15263) with regards to small remote-piloted-aerial-systems (RPAS); effective as of June 1st, 2019.
    \item Provided drone documentation is in compliance with Transport Canada's regulations regarding drone manufacturers; documentation include safety, pre-flight, radio, operation, post-flight, and emergency procedures for the drone.
    \item Drone is not over 25kg and will not fly over bystanders (anyone that is not actively participating in the flight operation) or within 30m of them.
    \item Drone will not fly within 3 nautical miles of any aerodrome.
    \item Drone will not fly over 400ft; and not fly in controlled airspaces: class A, class B, class C, class D, class E, and class F.
    \item Drone will only fly in fair weather conditions with adequate visibility for line-of-sight operations.
    \item Drone will not carry weapons and cannot be used or modified for weaponization.
\end{enumerate}

\subsection{Environmental Considerations}
\begin{enumerate}[{C.EN}.1:]
    \item 
\end{enumerate}

% Bibliography
\clearpage
\addcontentsline{toc}{section}{References}
\bibliographystyle{ieeetr}
\bibliography{references}

% Appendix (uncomment to enable appendix)
% \clearpage
% \appendix
% \section{Appendix name}\label{appendix:sample-appendix}
% Content here

\end{document}